\subsection{Web Application Testing}

The web app test file is designed to check the main user-facing behavior of the React application in a clear and controlled way. The tests render 
the app with mock data so we can verify behavior without needing live hardware or external services. The test file is organized under the test directory within the source directory,
along with the file \texttt{sampleWaterValues.js} that contains mock data. The test file uses the React Testing Library to render components and simulate user interactions,
and Vitest for mocks, assertions, and test organization.

\subsubsection{Basic UI Rendering}

This section verifies that the interface renders correctly in both empty and normal states. It checks the message shown when no reading exists, 
then confirms that water-level data, weight, and battery metadata are displayed when readings are available. It also checks that the displayed 
percentage matches the expected calculation.

\subsubsection{Calibration Actions and Edge Cases}

This section tests user interaction with calibration controls. The test suite clicks the calibrate empty, calibrate full, and reset buttons, 
then verifies that the expected calibration upsert operations are made. This section verifies calibration logic for both standard and edge 
conditions. The tests confirm that custom empty/full calibration values are used, and that percentage output is clamped correctly. Readings 
below empty are shown as 0\%, readings above full are shown as 100\%, and invalid ranges where full is less than or equal to empty also resolve to 0\%.

\subsubsection{Push Notification Behavior}

This section tests notification flows with mocked browser APIs. It includes unsupported-browser behavior, denied permission behavior, and successful 
permission behavior. In the successful case, the test verifies service worker registration and push subscription setup.

\subsubsection{API-Related Checks}

This section covers API behavior that is directly tested. Specifically, when notifications are enabled successfully, the app is verified to send a 
POST request to the token registration endpoint. Supabase reads and calibration writes are mocked so request behavior and state changes can be 
checked without live backend dependencies.
